\begin{longtable}[c]{| L{6cm} | L{6cm} |}
\caption{UART and LP UART Feature Comparison}
\label{tab:uart-feature-distinction} \\
    \hline
    \rowcolor{lightgray}
    \textbf{UART Feature} & \textbf{LP UART Feature}\\\hline
    \endfirsthead

    \multicolumn{2}{c}{\bfseries \tablename\ \thetable{} -- cont'd from previous page}\\\hline
    \rowcolor{lightgray}
    \textbf{UART Feature} & \textbf{LP\_UART Feature}\\\hline
    \endhead

    \multicolumn{2}{c}{\textbf{Cont'd on next page}} \\
    \endfoot
    \endlastfoot

\multicolumn{2}{|c|}{Programmable baud rate up to 5 MBaud} \\ \hline
128 x 8 bit RAM respectively for the TX channel and RX channel of a UART controller & 20 x 8-bit RAM, shared by the TX FIFO and RX FIFO of LP UART \\ \hline
\multicolumn{2}{|c|}{Full-duplex asynchronous communication} \\ \hline
\multicolumn{2}{|c|}{Data bits (5 to 8 bits)} \\ \hline
\multicolumn{2}{|c|}{Stop bits (1, 1.5, or 2 bits)} \\ \hline
\multicolumn{2}{|c|}{Parity bit} \\ \hline
\multicolumn{2}{|c|}{Special character AT\_CMD detection} \\ \hline
RS485 protocol & --- \\ \hline
IrDA protocol & --- \\ \hline
High-speed data communication using GDMA & --- \\ \hline
\multicolumn{2}{|c|}{Receive timeout} \\\hline
\multicolumn{2}{|c|}{UART as wakeup source} \\ \hline
\multicolumn{2}{|c|}{Software and hardware flow control} \\ \hline
\makecell[l]{Three prescalable clock sources:\\1. XTAL\_CLK\\ 2. RC\_FAST\_CLK\\ 3. PLL\_F80M\_CLK} & \makecell[l]{Three prescalable clock sources\\1. RC\_FAST\_CLK\\2. XTAL\_DIV\_CLK\\3. PLL\_F8M\_CLK} \\ \hline
\end{longtable}
